\documentclass[11pt,a4paper]{article}
\usepackage[margin=25mm]{geometry}
\usepackage{amsmath,amssymb,amsthm,hyperref}
\newtheorem{lemma}{Lemma}
\newcommand{\gen}[1]{\langle#1\rangle}
\title{Commuting interval factors and the crossing anchor}
\author{Human Compression project: bounded L044 proof component}
\date{}
\begin{document}\maketitle
Use the complete block-deficit setting: $G\leq A_n$ is a transitive
string C-group, $n\geq12$, every maximal omission is intransitive,
and a unique eligible $i$-edge joins $a\in A$ to $b\in B$, the two
$G_i$-orbits. Their two-block systems have block sizes $k,l$, respectively. Write $L=\gen{\rho_j:j<i}$
and $R=\gen{\rho_j:j>i}$. These commute, and $G_i=LR$.

\begin{lemma}[An anchor fixes the whole commuting orbit]
If every generator of $R$ fixes $a$, then $R$ fixes $A$ pointwise.
The symmetric assertion holds for $L$ and $b$.
\end{lemma}
\begin{proof}
For $l\in L$ and $r\in R$, commutation gives
$(a^l)^r=(a^r)^l=a^l$. Moreover $a^{lr}=a^l$, so the
$L$-orbit of $a$ is its whole $G_i$-orbit $A$.
Every element of $R$ therefore fixes every point of $A$.
\end{proof}

\begin{lemma}[No regular orbit away from its anchor]
Suppose $\alpha=0$, $\beta=1$, $\gamma=0$. The unique regular
four-point orbit $D$ of $M\cong C_2^2$ must contain the crossing
endpoint in its original $G_i$-orbit.
\end{lemma}
\begin{proof}
Suppose instead $D\subset A$ and $a\notin D$.
Both anchors then lie in paired $M$-orbits. Only the literal
neighbors $\rho_{i-1},\rho_{i+1}$ can exchange blocks: every
other generator fixes both anchors. The two independent
block-image generators are precisely those neighbors.
If they both exchange the blocks of either orbit, they act
identically on its anchor pair. Every other generator fixes
its opposite anchor, using commutation with one neighbor or,
for either adjacent exception, with the opposite neighbor.
The pair would be $G_i$-invariant, impossible.
Thus the actual block-image basis is pure. After reversal,
$\tau_A=\rho_{i-1}$ exchanges only the blocks of $A$, and
$\tau_B=\rho_{i+1}$ exchanges only the blocks of $B$.

On each paired $M$-orbit in $A$, projection to its two blocks
is an equivariant bijection. The element $\tau_B$ consequently
fixes that pair pointwise, in particular fixing $a$.
All other generators of $R$ fix $a$ by the unique-crossing
property. Lemma~1 therefore makes $R$ fix $A$ pointwise.
This is impossible: on the regular $M$-orbit $D$, the
nonidentity element $\tau_B$ has two transpositions.
\end{proof}

\begin{lemma}[Exact support separation for paired rows]
If $\alpha=\beta=0$, the literal pure block-image basis established
in the two-swap saturation note gives $L$ supported on $A$ and
$R$ supported on $B$. Their actions there are faithful transitive
string C-markings with their original generator orders.
\end{lemma}
\begin{proof}
The right neighbor fixes $A$ pointwise, and all other generators
of $R$ fix $a$. Lemma~1 makes all of $R$ trivial on $A$.
The same argument makes $L$ trivial on $B$.
Faithfulness follows from these being the only two original
orbits. Transitivity follows from $G_i=LR$; string relations,
independence and IP are inherited by the unchanged intervals.
\end{proof}

\begin{lemma}[The paired-row rank bound]
If $\alpha=\beta=0$, then
$r\leq k+l-1=(n-2)/2$.
\end{lemma}
\begin{proof}
The marked pure-basis and parity arguments in the two-swap note give
even $k,l\geq4$. By Lemma~3, the left interval is supported on $A$
and transitive there; the right interval is supported on $B$ and
transitive there. Let their ranks be $d_A,d_B$, so
$r=d_A+d_B+1$.

The left interval ends with $x=\rho_{i-2}$ and
$\tau_A=\rho_{i-1}$, both present. All its earlier generators
commute with $\tau_A$ and fix $a$, so act faithfully and diagonally
on the other $k-1$ row pairs. Their independent length is at most
$k-2$ by Whiston. Hence $d_A\leq k$. Suppose equality holds.
The earlier family then generates the full $S_{k-1}$, with a
maximal independent involutory string marking.

For $k\geq6$, the accepted L047 classification makes this the
ordinary Coxeter path: $k-1$ is odd, so the outer degree-six
exception is absent. Its near-end omission is the natural
$S_{k-2}$ fixing the anchor row and one additional row $v$.
The generator $x$ centralizes that omission, preserves the two
original blocks and fixes $a$. It is consequently the identity
on the block containing $a$, and on the other block it is the
transposition exchanging $a'$ with $v'$, since it must move
$a'$ to another row. Its support is contained in $A$ by Lemma~3.
It is therefore odd, impossible.

For $k=4$, denote the far and near tail generators by $p,q$.
Choose row labels $0,1,2,3$ and $0',1',2',3'$ with $a=0$, so
\[
 p=(1\ 2)(1'\ 2'),\quad q=(2\ 3)(2'\ 3'),\quad
 \tau_A=(0\ 0')(1\ 1')(2\ 2')(3\ 3').
\]
Since $x$ centralizes $p$, fixes $0$, moves $0'$ to another row,
and is even, its only possibilities are
\[
               x=(1\ 2)(0'\ 3')
                 \quad\hbox{or}\quad(1'\ 2')(0'\ 3').
\]
In either case $\gen{q,x,\tau_A}$ is transitive on all eight
points. This follows directly from its permutation graph:
contract the four $\tau_A$ pairs to row vertices. The $q$ edges
join rows $2,3$, and the $x$ edges join $0,3$ and $1,2$.
These three row edges form a spanning path, and each contracted
pair is itself joined by $\tau_A$. The original graph is thus
connected, proving transitivity without an order computation.

Let $j$ be the index of $p$. The global omission $G_j$ contains
this transitive subgroup on $A$, the unchanged entire right
interval transitive on $B$, and $\rho_i$ joining $a$ to $b$.
It is therefore transitive on $A\cup B$, contradicting the
standing hypothesis on every maximal omission. This excludes
$d_A=k$ also at $k=4$.

We have proved $d_A\leq k-1$. The symmetric proof gives
$d_B\leq l-1$, whence $r\leq k+l-1$ as asserted.
\end{proof}

\paragraph{Scope and attribution.}
The source problem and block-image framework are from Cameron,
Fernandes, Leemans and Mixer, \emph{Highest rank of a polytope
for $A_n$}, arXiv:1605.09173v1, Section~5.1.
The elementary commuting-factor argument supplies a short new
proof in that framework. The complete accepted block-deficit
and new frozen two-swap components remain paid dependencies.
Together with the anchored regular-orbit proof, Lemma~2 closes
$\beta=1$ for two blocks on each original orbit. Lemma~4 closes
the remaining $\alpha=\beta=0$ paired-row case. Combined with
the earlier $\alpha=1$ endpoint results and the exact deficit
formula, these prove the rank bound when both original orbits
have two blocks. Mixed block-count and other source branches remain.
Full-source recovery,
new computation, new formalization and new-question credit are
not claimed.
\end{document}
